\chapter{Page Explorer : Le Catalogue de Données}
\label{chap:explorer}

La page \textbf{Explorer} (\texttt{explorer.html}) constitue le point d'entrée principal et la tour de contrôle de la WebPlatform. C'est ici que les chercheurs de l'IRIBHM parcourent, filtrent et sélectionnent les données biologiques de souris en cours d'étude. Ce chapitre détaille l'ergonomie, la structure mémoire de l'index client et les algorithmes de prétraitement automatisés.

\section{Interface et Ergonomie Utilisateur}

L'interface graphique est orchestrée par le script \texttt{js/pages/explorer.js} qui manipule dynamiquement le DOM à partir de l'état mémoire du catalogue. Elle est divisée en modules hautement interactifs.

\subsection{Barre de Recherche et Mécanisme de Debounce}

La recherche textuelle permet d'explorer en temps réel les champs textuels du catalogue (nom, description, canaux, marqueurs moléculaires, région anatomique). Pour éviter de recalculer le rendu de la grille à chaque frappe de touche (ce qui nuirait à la fluidité de saisie), l'événement clavier est filtré par un mécanisme de \textbf{Debounce} temporel :

\begin{equation}
    T_{\text{last\_keypress}} - T_{\text{last\_execution}} > \tau_{\text{debounce}} \quad (\text{où } \tau_{\text{debounce}} = 300\,\text{ms})
\end{equation}

Si une nouvelle touche est pressée avant l'expiration de l'intervalle $\tau_{\text{debounce}}$, le timer précédent est annulé et réinitialisé. Cette temporisation évite les saccades d'interface en limitant les appels à la fonction de filtrage.

\subsection{Système de Filtrage Multidimensionnel}

Les filtres à facettes permettent de croiser plusieurs critères :
\begin{itemize}
    \item \textbf{Type d'Acquisition} :
        \begin{itemize}
            \item \textit{Fixed Imaging} : Données spatiales tridimensionnelles de haute résolution.
            \item \textit{Live Imaging} : Séries temporelles 4D d'embryons en développement.
            \item \textit{Cell Tracking} : Enregistrements possédant des données de trajectoires cellulaires quantitatives.
        \end{itemize}
    \item \textbf{Stade Embryonnaire (Stage)} : Les stades de développement (ex : E10.5, E11.5) sont extraits de manière unique et dynamique à partir de tous les éléments du catalogue. L'interface construit à la volée une liste déroulante d'options.
\end{itemize}

\subsection{Cartographie d'une Carte de Jeu de Données}

Chaque jeu de données est matérialisé par un composant carte (Card) affichant :
\begin{itemize}
    \item \textbf{Vignette (Thumbnail)} : Image de synthèse représentative du volume.
    \item \textbf{Badges de statut} : `Raw` (présence du fichier original volumineux), `Web` (briques prêtes pour l'affichage en ligne), `Tracking` (fichier de trajectoires associé).
    \item \textbf{Visualisation des canaux} : Des indicateurs colorés représentent les canaux fluorescents présents (ex: DAPI en bleu, GFP en vert, Pecam1 en rouge), permettant au chercheur de connaître immédiatement les protéines cibles marquées dans l'échantillon.
\end{itemize}

\section{Architecture Interne du Catalogue Client (\texttt{catalog.js})}

L'application web fonctionne selon un modèle local sans base de données SQL active côté serveur. Cela garantit une portabilité totale de la plateforme (elle peut être exécutée directement depuis un disque local ou une clé USB).

\subsection{Chargement et Cycle de Vie}
Au démarrage, \texttt{js/core/catalog.js} effectue une requête asynchrone \texttt{fetch} pour charger le fichier JSON maître situé à l'adresse \texttt{DATA\_WEB/catalog.json}. Cet objet JSON est stocké en mémoire dans un tableau privé \texttt{\_datasets} et sert de source unique de vérité.

\subsection{Structure du Schéma du Catalogue}
Voici un exemple d'enregistrement valide pour un volume confocal complexe dans \texttt{catalog.json} :

\begin{lstlisting}[language=JavaScript, caption=Extrait représentatif du schéma catalog.json]
{
  "id": "2026_E10.5_LimbBud_GFP",
  "name": "Bourgeon de membre - E10.5",
  "description": "Acquisition haute résolution du système vasculaire du bourgeon de membre",
  "type": "fixed",
  "date": "2026-05-28",
  "stage": "E10.5",
  "physicalSizeUm": {
    "x": 0.35,
    "y": 0.35,
    "z": 1.50
  },
  "dimensions": {
    "width": 1024,
    "height": 1024,
    "depth": 150
  },
  "channels": [
    { "name": "DAPI", "excitation": 405, "color": "#0000ff" },
    { "name": "Pecam1", "excitation": 561, "color": "#ff0000" }
  ],
  "path": "fixed/2026_E10.5_LimbBud_GFP",
  "thumbnail": "DATA_WEB/fixed/2026_E10.5_LimbBud_GFP/thumbnail.webp"
}
\end{lstlisting}

\subsection{Algorithme de Filtrage Client}
Le filtrage combine l'intersection des facettes actives et une recherche textuelle insensible à la casse.

\begin{lstlisting}[language=Java, caption=Algorithme complet de filtrage et recherche]
function getFilteredDatasets(criteria) {
    return _datasets.filter(dataset => {
        // 1. Filtrage par type
        if (criteria.type && criteria.type !== 'all' && dataset.type !== criteria.type) {
            return false;
        }
        // 2. Filtrage par stade embryonnaire
        if (criteria.stage && criteria.stage !== 'all' && dataset.stage !== criteria.stage) {
            return false;
        }
        // 3. Recherche textuelle multi-champs
        if (criteria.searchQuery) {
            const query = criteria.searchQuery.toLowerCase().trim();
            const fieldsToSearch = [
                dataset.name,
                dataset.description,
                dataset.stage,
                ...dataset.channels.map(c => c.name)
            ];
            const match = fieldsToSearch.some(field => 
                field && field.toLowerCase().includes(query)
            );
            if (!match) return false;
        }
        return true;
    });
}
\end{lstlisting}

\section{Génération du Catalogue Automatisée par Python}

Le fichier maître \texttt{catalog.json} n'est pas écrit à la main. Il est construit à chaque modification du stockage par le script Python \texttt{preprocess/generate\_catalog.py}.

\subsection{Algorithme de Projection d'Intensité Maximale (MIP) pour Vignettes}
Pour générer les thumbnails, le preprocessor n'effectue pas de downsampling destructif global du volume. À la place, il calcule une \textbf{Projection d'Intensité Maximale (MIP)} le long de l'axe optique (Z) du microscope. 
Mathématiquement, pour un volume discret tridimensionnel $V(x, y, z)$ de dimensions $N_x \times N_y \times N_z$, le pixel $I_{\text{MIP}}(x, y)$ de l'image 2D résultante est calculé par la formule :

\begin{equation}
    I_{\text{MIP}}(x, y) = \max_{z \in [0, N_z-1]} V(x, y, z)
\end{equation}

Cette opération est réalisée de manière optimisée sous Python en tirant parti du calcul vectorisé de NumPy :
\begin{lstlisting}[language=Python, caption=Calcul du MIP et conversion WebP via PIL]
import numpy as np
from PIL import Image

def generate_mip_thumbnail(volume_3d, output_path):
    # volume_3d est un array NumPy de forme (Z, Y, X)
    mip_array = np.max(volume_3d, axis=0) # Projection sur l'axe Z
    
    # Normalisation des intensites sur 8 bits
    mip_normalized = ((mip_array - mip_array.min()) / (mip_array.max() - mip_array.min()) * 255.0).astype(np.uint8)
    
    # Sauvegarde au format WebP compresse (sans perte d'information visuelle)
    img = Image.fromarray(mip_normalized)
    img.save(output_path, "WEBP", quality=85)
\end{lstlisting}

Cette projection garantit que toutes les structures fluorescentes brillantes (comme les vaisseaux capillaires ou les noyaux) restent parfaitement visibles sur la vignette 2D, même si elles se situent en profondeur dans le tissu.

\subsection{Validation de la Calibration Physique}
Le preprocessor calcule rigoureusement le volume physique en micromètres cubes ($V_{\text{physical}}$) pour s'assurer qu'aucune anomalie d'échelle ne perturbe le visualiseur :

\begin{equation}
    V_{\text{physical}} = (N_x \cdot s_x) \times (N_y \cdot s_y) \times (N_z \cdot s_z)
\end{equation}

où $s_x, s_y, s_z$ sont les résolutions spatiales par pixel (voxel sizes) extraites des métadonnées originales (comme les métadonnées TIFF ou XML d'acquisition Bio-Formats). Si le script Python détecte que les valeurs de résolutions physiques ne concordent pas avec les rapports d'échelles standards (ex : ratio d'aspect disproportionné en Z), il injecte un avertissement dans le fichier final pour que l'utilisateur en soit notifié lors de l'analyse 3D.
